\section{实数的无穷和与无穷乘积}

由于 \(\mathbf{R}\) 的每一点都有可数的基本邻域系（§ 1，第 4 节，命题 3 的推论），由此推出，有限实数的族 \((x_{t})\) 在 \(\mathbf{R}\) 中可和只当指标 \(t\) （满足 \(x_{t} \neq 0\) ）组成的集合是可数的（第三章，§ 5，第 2 节，命题 1 的推论）。于是 \(\mathbf{R}\) 中可和族的研究本质上归结为可和序列的研究。然而，以后我们不得不考虑不可数的有限实数族 \((x_{t})\) ，其各项是参数 \(t\) 的函数；可能发生这样的情形：对一切 t 该族可和，但由指标 \(\iota\) （满足 \(x_{t} \neq 0\) ）组成的（可数）集合依赖于 t。由于这个原因，下面我们对指标集的势不作任何假设。

